\chapter{There are infinitely many primes}
\label{ch:infinitude}

\section{The source's argument, stated so that it is true}

The source argues by contradiction from ``there are only finitely many primes
$p_1 \dots p_n$'', forms $m = \prod_i p_i$, and factors $m + 1$. Formalising
this needs two things the source leaves implicit.

\begin{itemize}
  \item \textbf{$1 < m + 1$ must be discharged} before
    Corollary~\ref{cor:exists_prime_dvd} applies, i.e.\ $m \neq 0$. This is
    \emph{not} automatic for an arbitrary list: a list that merely
    \emph{contains} every prime could also contain $0$, and then $m = 0$ and
    $m + 1 = 1$ has no prime divisor at all. We therefore hypothesise that the
    list contains \emph{only} primes as well as \emph{all} primes. Since every
    prime is $> 1 > 0$, the product of such a list is positive.
  \item \textbf{The contradiction is $q \mid 1$ against $1 < q$}, so it is again
    the $1 < p$ conjunct of Definition~\ref{def:is_prime} that carries the
    argument. Without it the proof has no contradiction to reach.
\end{itemize}

\begin{theorem}[no list contains exactly the primes]
  \leanok
  \label{thm:not_exists_prime_list}
  \lean{EuclidPrimes.not_exists_isPrime_list}
  \uses{def:is_prime, cor:exists_prime_dvd}
  % SOURCE: [euclid-notes], the final `theorem:` (read from references/euclid-infinitude-of-primes.md)
  % SOURCE QUOTE: "theorem: There are inifinitely many primes."
  \textit{Source: euclid-infinitude-of-primes.md, final \texttt{theorem:}.}
  There is no finite list $\ell$ of natural numbers such that
  \[
    (\forall p \in \ell,\ \Prm{p}) \quad\text{and}\quad (\forall p,\ \Prm{p} \implies p \in \ell).
  \]
\end{theorem}
% SOURCE QUOTE PROOF: "Suppose there are only finitely many primes p1...pn. let m = (\prod_{i=1}^{n} p_i). Then n admits a prime factorization
%
%   m + 1 = \Product_{j=1}^{k} p_j ^ e_j
%
% . Then for all p_j, p_j | m + 1, p_j | m => m_j | 1, contradiction."
\begin{proof}
  \uses{cor:exists_prime_dvd}
  \leanok
  Suppose such an $\ell$ exists, and set $m := \prod \ell$.

  \paragraph{Step 1: $0 < m$, hence $1 < m + 1$.}
  Every member of $\ell$ is prime by the first hypothesis, hence $> 1$ and in
  particular $> 0$. A product of positive naturals is positive, so $0 < m$ and
  therefore $1 < m + 1$.

  \paragraph{Step 2: $m + 1$ has a prime divisor.}
  By Corollary~\ref{cor:exists_prime_dvd} applied to $m + 1$ --- legitimate
  exactly because of Step 1 --- there is $q$ with $\Prm{q}$ and $q \mid m + 1$.

  \paragraph{Step 3: $q$ also divides $m$.}
  By the second hypothesis, $\Prm{q}$ gives $q \in \ell$, and any member of a
  list divides that list's product, so $q \mid m$.

  \paragraph{Step 4: contradiction.}
  $q$ divides both $m + 1$ and $m$, hence divides their difference
  $(m+1) - m = 1$. A divisor of $1$ is at most $1$, contradicting $1 < q$ from
  $\Prm{q}$.
\end{proof}

\section{The theorem in its usual forms}

\begin{corollary}[the set of primes is infinite]
  \leanok
  \label{cor:infinite_setOf_prime}
  \lean{EuclidPrimes.infinite_setOf_isPrime}
  \uses{def:is_prime, thm:not_exists_prime_list}
  The set $\{\, p \in \NN \mid \Prm{p} \,\}$ is infinite.
\end{corollary}
\begin{proof}
  \uses{thm:not_exists_prime_list}
  \leanok
  Suppose it were finite. A finite set of natural numbers can be enumerated as a
  finite list $\ell$ whose members are exactly its elements. That $\ell$ then
  satisfies both hypotheses of Theorem~\ref{thm:not_exists_prime_list}: its
  members are primes, and every prime, being an element of the set, occurs in
  it. This contradicts the theorem.
\end{proof}

\begin{corollary}[primes are unbounded]
  \leanok
  \label{cor:exists_prime_gt}
  \lean{EuclidPrimes.exists_isPrime_gt}
  \uses{def:is_prime, cor:infinite_setOf_prime}
  For every $n \in \NN$ there exists $p$ with $\Prm{p}$ and $n < p$.
\end{corollary}
\begin{proof}
  \uses{cor:infinite_setOf_prime}
  \leanok
  An infinite set of natural numbers contains an element greater than any given
  $n$; apply this to Corollary~\ref{cor:infinite_setOf_prime}.
\end{proof}
